\documentclass[11pt,a4paper]{article}
\usepackage[utf8]{inputenc}
\usepackage[T1]{fontenc}
\usepackage[spanish,es-minimal,es-noshorthands]{babel}
\AtBeginDocument{\renewcommand{\sin}{\operatorname{sen}}\renewcommand{\arcsin}{\operatorname{arcsen}}}
\usepackage{lmodern}
\usepackage[sfdefault,lining]{FiraSans}
\usepackage[a4paper,top=2.3cm,bottom=2.2cm,left=1.8cm,right=1.8cm,headheight=1.3cm,headsep=0.4cm,footskip=1cm]{geometry}
\usepackage{amsmath,amssymb}
\usepackage{icomma}
\usepackage{graphicx}
\usepackage[table]{xcolor}
\usepackage{tikz}
\usepackage{fontawesome5}
\usepackage{ccicons}
\usepackage{fancyhdr}
\usepackage{lastpage}
\usepackage{tasks}
\usepackage{needspace}
\usepackage[most]{tcolorbox}
\usepackage{array}
\usepackage[hidelinks]{hyperref}
\AtBeginDocument{\hypersetup{pdftitle={Soluciones de la ficha de la unidad 5: Trigonometría. 1º Bach.},pdfsubject={Matemáticas I 1º Bach.},pdfkeywords={trigonometría, radianes, razones trigonométricas, circunferencia goniométrica, resolución de triángulos, teorema del seno, teorema del coseno, fórmulas trigonométricas, ecuaciones trigonométricas, funciones trigonométricas}}}

\definecolor{navy}{HTML}{1B2748}
\definecolor{accent}{HTML}{E8603A}
\definecolor{teal}{HTML}{1F9E90}
\definecolor{soft}{HTML}{EEF1F8}
\definecolor{muted}{HTML}{5D6685}

\setlength{\parindent}{0pt}
\setlength{\parskip}{3pt}
\settasks{label=\alph*), label-format=\bfseries\color{navy}, label-width=1.3em, item-indent=1.8em, column-sep=0.6em, after-item-skip=5pt, before-skip=3pt, after-skip=4pt}

% ---------- capçalera i peu ----------
\newcommand{\llicencia}{\href{https://creativecommons.org/licenses/by-nc-sa/4.0/deed.es}{\ccbyncsa}\ \ Licencia Creative Commons Atribución-NoComercial-CompartirIgual 4.0 Internacional (CC BY-NC-SA 4.0)}
\pagestyle{fancy}
\fancyhf{}
\fancyhead[L]{\raisebox{-0.35\height}{\includegraphics[height=0.95cm]{logo-gva.png}}}
\fancyhead[C]{\small\color{muted}Soluciones de la ficha · Unidad 5: Trigonometría · 1\textsuperscript{o} Bach.}
\fancyhead[R]{\raisebox{-0.35\height}{\includegraphics[height=0.95cm]{logo-vermellar.png}}}
\fancyfoot[L]{\scriptsize\color{muted}\llicencia}
\fancyfoot[R]{\small\color{muted}\thepage}
\renewcommand{\headrulewidth}{0.4pt}
\renewcommand{\headrule}{\color{muted!40}\hrule height \headrulewidth}
\fancypagestyle{primera}{\fancyhf{}\renewcommand{\headrulewidth}{0pt}
  \fancyfoot[L]{\scriptsize\color{muted}\llicencia}
  \fancyfoot[R]{\small\color{muted}\thepage}}

% ---------- nivells i exercicis ----------
\newcommand{\estrelles}[1]{{\color{accent}\ifcase#1\or\faStar\faStar[regular]\faStar[regular]\or\faStar\faStar\faStar[regular]\or\faStar\faStar\faStar\fi}}
\newcounter{exer}
\newcommand{\ex}[1]{\par\needspace{4\baselineskip}\medskip\refstepcounter{exer}\noindent{\color{navy}\bfseries\theexer.}\ \ }

\newcommand{\seccio}[2]{\par\needspace{16\baselineskip}\bigskip
  \begin{tcolorbox}[colback=navy,colframe=navy,coltext=white,arc=2mm,boxrule=0pt,left=3mm,right=3mm,top=1.2mm,bottom=1.2mm]
  {\large\bfseries #1\quad #2}
  \end{tcolorbox}\vspace{-2pt}}

\newcommand{\exf}[1]{\par\medskip\refstepcounter{exer}\noindent{\color{navy}\bfseries\theexer.}\ \ }
\usetikzlibrary{calc}
\newcolumntype{C}[1]{>{\centering\arraybackslash}p{#1}}
% quadre de recordatori (només a la fitxa d'exercicis)
\newcommand{\recorda}[1]{\par\smallskip\begin{tcolorbox}[colback=accent!7,colframe=accent!55,arc=2mm,boxrule=0.5pt,left=2.5mm,right=2.5mm,top=1mm,bottom=1mm]\small{\color{accent}\faLightbulb[regular]}\ \textbf{Recuerda:}\ #1\end{tcolorbox}}
\newcommand{\res}[1]{\boxed{#1}}\newcommand{\rest}[1]{\fbox{#1}}

\def\fmneed{11}% que el títol d'una secció no quede sol al final del full
% =====================================================================
% Capçalera FentMates de les fitxes: un color per a cada curs.
% Els colors són els de la pàgina inicial (LEVELS a assets/js/materials.js);
% el logo ve de _eines/logo/pdf/ (python3 _eines/logo/genera.py).
%
% Ús, al final del preàmbul de la unitat:
%   % =====================================================================
% Capçalera FentMates de les fitxes: un color per a cada curs.
% Els colors són els de la pàgina inicial (LEVELS a assets/js/materials.js);
% el logo ve de _eines/logo/pdf/ (python3 _eines/logo/genera.py).
%
% Ús, al final del preàmbul de la unitat:
%   % =====================================================================
% Capçalera FentMates de les fitxes: un color per a cada curs.
% Els colors són els de la pàgina inicial (LEVELS a assets/js/materials.js);
% el logo ve de _eines/logo/pdf/ (python3 _eines/logo/genera.py).
%
% Ús, al final del preàmbul de la unitat:
%   \input{../comu/fentmates.tex}
%   \fmcurs{1eso}
% i, al començament del document:
%   \capcalera{Matemàtiques}{1r ESO}{Fitxa d'exercicis · Unitat 2}{Divisibilitat}
% Opcional abans de \input: \def\fmneed{14} (espai mínim abans de cada secció).
% =====================================================================
\makeatletter
\def\fm@logodir{../../logo/pdf}
\providecommand{\fmneed}{10}
\def\fm@cursnom{}

% clau del curs, color i color del text damunt del color
\newcommand{\fm@curs}[3]{\@namedef{fm@color@#1}{#2}\@namedef{fm@text@#1}{#3}}
\fm@curs{1eso}{E6B84A}{navy}
\fm@curs{2eso}{3B82F6}{white}
\fm@curs{3eso}{FF7A50}{white}
\fm@curs{4eso-a}{A78BFA}{white}
\fm@curs{4eso-b}{6D28D9}{white}
\fm@curs{1batx-mi}{9B1C1C}{white}
\fm@curs{1batx-ccss}{EF6B6B}{white}
\fm@curs{2batx-mii}{1E6B47}{white}
\fm@curs{2batx-ccss}{2BB3A3}{white}

\newcommand{\fmcurs}[1]{%
  \@ifundefined{fm@color@#1}{\PackageError{fentmates}{Curs desconegut: #1}{}}{}%
  \edef\fm@tmp{\noexpand\definecolor{curs}{HTML}{\@nameuse{fm@color@#1}}%
               \noexpand\colorlet{curstext}{\@nameuse{fm@text@#1}}}\fm@tmp
  \colorlet{cursfosc}{curs!72!black}%
  \def\fm@logo{\fm@logodir/vertical-#1.pdf}%
  \def\fm@marca{\fm@logodir/marca-#1.pdf}}

% ---------- capçalera de la primera pàgina ----------
% #1 matèria (petit), #2 curs, #3 tipus de document i unitat, #4 títol
\newcommand{\capcalera}[4]{%
  \gdef\fm@cursnom{#2}%
  \par\noindent
  \begin{tcolorbox}[colback=curs,colframe=curs,boxrule=0pt,arc=2.2mm,width=3.5cm,height=2.15cm,
      valign=center,left=3mm,right=2mm,top=1mm,bottom=1mm,nobeforeafter,box align=center]
    \color{curstext}{\footnotesize #1}\par\vspace{1pt}{\fontsize{19}{21}\selectfont\bfseries #2\par}%
  \end{tcolorbox}\hspace{5mm}%
  \begin{minipage}[c]{\dimexpr\linewidth-3.5cm-5mm-2.9cm\relax}
    {\color{cursfosc}\small\bfseries\MakeUppercase{#3}}\par\vspace{2pt}
    {\color{navy}\fontsize{20}{23}\selectfont\bfseries #4\par}
  \end{minipage}\hfill
  \begin{minipage}[c]{2.6cm}\raggedleft\includegraphics[width=2.6cm]{\fm@logo}\end{minipage}\par
  \vspace{2pt}}

% ---------- barres de secció del color del curs ----------
\renewcommand{\seccio}[2]{\par\needspace{\fmneed\baselineskip}\bigskip
  \begin{tcolorbox}[colback=curs,colframe=curs,coltext=curstext,arc=2mm,boxrule=0pt,left=3mm,right=3mm,top=1.2mm,bottom=1.2mm]
  {\large\bfseries #1\quad #2}
  \end{tcolorbox}\vspace{-2pt}}

% ---------- estrelles de dificultat del color del curs (a Batxillerat no se'n posen) ----------
\@ifundefined{estrelles}{}{%
  \renewcommand{\estrelles}[1]{{\color{curs}\ifcase#1\or\faStar\faStar[regular]\faStar[regular]\or\faStar\faStar\faStar[regular]\or\faStar\faStar\faStar\fi}}}

% ---------- capçalera i peu de les pàgines ----------
% primera pàgina: logos institucionals al peu
\fancypagestyle{primera}{\fancyhf{}\renewcommand{\headrulewidth}{0pt}
  \fancyfoot[L]{\raisebox{-0.32\height}{\includegraphics[height=0.78cm]{logo-gva.png}}\hspace{5mm}%
    \raisebox{-0.32\height}{\includegraphics[height=0.78cm]{logo-vermellar.png}}}
  \fancyfoot[R]{\scriptsize\color{muted}\href{https://creativecommons.org/licenses/by-nc-sa/4.0/deed.ca}{\ccbyncsa}\ \ CC BY-NC-SA 4.0 · fentmates.org\qquad\small\thepage}}
% la resta: el símbol a l'esquerra i el curs a la dreta, amb el títol al centre com abans
\fancyhead[L]{\raisebox{-0.3\height}{\includegraphics[height=0.62cm]{\fm@marca}}}
\fancyhead[R]{\small\bfseries\color{cursfosc}\fm@cursnom}
\renewcommand{\headrule}{\color{curs!45}\hrule height \headrulewidth}
\makeatother

%   \fmcurs{1eso}
% i, al començament del document:
%   \capcalera{Matemàtiques}{1r ESO}{Fitxa d'exercicis · Unitat 2}{Divisibilitat}
% Opcional abans de \input: \def\fmneed{14} (espai mínim abans de cada secció).
% =====================================================================
\makeatletter
\def\fm@logodir{../../logo/pdf}
\providecommand{\fmneed}{10}
\def\fm@cursnom{}

% clau del curs, color i color del text damunt del color
\newcommand{\fm@curs}[3]{\@namedef{fm@color@#1}{#2}\@namedef{fm@text@#1}{#3}}
\fm@curs{1eso}{E6B84A}{navy}
\fm@curs{2eso}{3B82F6}{white}
\fm@curs{3eso}{FF7A50}{white}
\fm@curs{4eso-a}{A78BFA}{white}
\fm@curs{4eso-b}{6D28D9}{white}
\fm@curs{1batx-mi}{9B1C1C}{white}
\fm@curs{1batx-ccss}{EF6B6B}{white}
\fm@curs{2batx-mii}{1E6B47}{white}
\fm@curs{2batx-ccss}{2BB3A3}{white}

\newcommand{\fmcurs}[1]{%
  \@ifundefined{fm@color@#1}{\PackageError{fentmates}{Curs desconegut: #1}{}}{}%
  \edef\fm@tmp{\noexpand\definecolor{curs}{HTML}{\@nameuse{fm@color@#1}}%
               \noexpand\colorlet{curstext}{\@nameuse{fm@text@#1}}}\fm@tmp
  \colorlet{cursfosc}{curs!72!black}%
  \def\fm@logo{\fm@logodir/vertical-#1.pdf}%
  \def\fm@marca{\fm@logodir/marca-#1.pdf}}

% ---------- capçalera de la primera pàgina ----------
% #1 matèria (petit), #2 curs, #3 tipus de document i unitat, #4 títol
\newcommand{\capcalera}[4]{%
  \gdef\fm@cursnom{#2}%
  \par\noindent
  \begin{tcolorbox}[colback=curs,colframe=curs,boxrule=0pt,arc=2.2mm,width=3.5cm,height=2.15cm,
      valign=center,left=3mm,right=2mm,top=1mm,bottom=1mm,nobeforeafter,box align=center]
    \color{curstext}{\footnotesize #1}\par\vspace{1pt}{\fontsize{19}{21}\selectfont\bfseries #2\par}%
  \end{tcolorbox}\hspace{5mm}%
  \begin{minipage}[c]{\dimexpr\linewidth-3.5cm-5mm-2.9cm\relax}
    {\color{cursfosc}\small\bfseries\MakeUppercase{#3}}\par\vspace{2pt}
    {\color{navy}\fontsize{20}{23}\selectfont\bfseries #4\par}
  \end{minipage}\hfill
  \begin{minipage}[c]{2.6cm}\raggedleft\includegraphics[width=2.6cm]{\fm@logo}\end{minipage}\par
  \vspace{2pt}}

% ---------- barres de secció del color del curs ----------
\renewcommand{\seccio}[2]{\par\needspace{\fmneed\baselineskip}\bigskip
  \begin{tcolorbox}[colback=curs,colframe=curs,coltext=curstext,arc=2mm,boxrule=0pt,left=3mm,right=3mm,top=1.2mm,bottom=1.2mm]
  {\large\bfseries #1\quad #2}
  \end{tcolorbox}\vspace{-2pt}}

% ---------- estrelles de dificultat del color del curs (a Batxillerat no se'n posen) ----------
\@ifundefined{estrelles}{}{%
  \renewcommand{\estrelles}[1]{{\color{curs}\ifcase#1\or\faStar\faStar[regular]\faStar[regular]\or\faStar\faStar\faStar[regular]\or\faStar\faStar\faStar\fi}}}

% ---------- capçalera i peu de les pàgines ----------
% primera pàgina: logos institucionals al peu
\fancypagestyle{primera}{\fancyhf{}\renewcommand{\headrulewidth}{0pt}
  \fancyfoot[L]{\raisebox{-0.32\height}{\includegraphics[height=0.78cm]{logo-gva.png}}\hspace{5mm}%
    \raisebox{-0.32\height}{\includegraphics[height=0.78cm]{logo-vermellar.png}}}
  \fancyfoot[R]{\scriptsize\color{muted}\href{https://creativecommons.org/licenses/by-nc-sa/4.0/deed.ca}{\ccbyncsa}\ \ CC BY-NC-SA 4.0 · fentmates.org\qquad\small\thepage}}
% la resta: el símbol a l'esquerra i el curs a la dreta, amb el títol al centre com abans
\fancyhead[L]{\raisebox{-0.3\height}{\includegraphics[height=0.62cm]{\fm@marca}}}
\fancyhead[R]{\small\bfseries\color{cursfosc}\fm@cursnom}
\renewcommand{\headrule}{\color{curs!45}\hrule height \headrulewidth}
\makeatother

%   \fmcurs{1eso}
% i, al començament del document:
%   \capcalera{Matemàtiques}{1r ESO}{Fitxa d'exercicis · Unitat 2}{Divisibilitat}
% Opcional abans de \input: \def\fmneed{14} (espai mínim abans de cada secció).
% =====================================================================
\makeatletter
\def\fm@logodir{../../logo/pdf}
\providecommand{\fmneed}{10}
\def\fm@cursnom{}

% clau del curs, color i color del text damunt del color
\newcommand{\fm@curs}[3]{\@namedef{fm@color@#1}{#2}\@namedef{fm@text@#1}{#3}}
\fm@curs{1eso}{E6B84A}{navy}
\fm@curs{2eso}{3B82F6}{white}
\fm@curs{3eso}{FF7A50}{white}
\fm@curs{4eso-a}{A78BFA}{white}
\fm@curs{4eso-b}{6D28D9}{white}
\fm@curs{1batx-mi}{9B1C1C}{white}
\fm@curs{1batx-ccss}{EF6B6B}{white}
\fm@curs{2batx-mii}{1E6B47}{white}
\fm@curs{2batx-ccss}{2BB3A3}{white}

\newcommand{\fmcurs}[1]{%
  \@ifundefined{fm@color@#1}{\PackageError{fentmates}{Curs desconegut: #1}{}}{}%
  \edef\fm@tmp{\noexpand\definecolor{curs}{HTML}{\@nameuse{fm@color@#1}}%
               \noexpand\colorlet{curstext}{\@nameuse{fm@text@#1}}}\fm@tmp
  \colorlet{cursfosc}{curs!72!black}%
  \def\fm@logo{\fm@logodir/vertical-#1.pdf}%
  \def\fm@marca{\fm@logodir/marca-#1.pdf}}

% ---------- capçalera de la primera pàgina ----------
% #1 matèria (petit), #2 curs, #3 tipus de document i unitat, #4 títol
\newcommand{\capcalera}[4]{%
  \gdef\fm@cursnom{#2}%
  \par\noindent
  \begin{tcolorbox}[colback=curs,colframe=curs,boxrule=0pt,arc=2.2mm,width=3.5cm,height=2.15cm,
      valign=center,left=3mm,right=2mm,top=1mm,bottom=1mm,nobeforeafter,box align=center]
    \color{curstext}{\footnotesize #1}\par\vspace{1pt}{\fontsize{19}{21}\selectfont\bfseries #2\par}%
  \end{tcolorbox}\hspace{5mm}%
  \begin{minipage}[c]{\dimexpr\linewidth-3.5cm-5mm-2.9cm\relax}
    {\color{cursfosc}\small\bfseries\MakeUppercase{#3}}\par\vspace{2pt}
    {\color{navy}\fontsize{20}{23}\selectfont\bfseries #4\par}
  \end{minipage}\hfill
  \begin{minipage}[c]{2.6cm}\raggedleft\includegraphics[width=2.6cm]{\fm@logo}\end{minipage}\par
  \vspace{2pt}}

% ---------- barres de secció del color del curs ----------
\renewcommand{\seccio}[2]{\par\needspace{\fmneed\baselineskip}\bigskip
  \begin{tcolorbox}[colback=curs,colframe=curs,coltext=curstext,arc=2mm,boxrule=0pt,left=3mm,right=3mm,top=1.2mm,bottom=1.2mm]
  {\large\bfseries #1\quad #2}
  \end{tcolorbox}\vspace{-2pt}}

% ---------- estrelles de dificultat del color del curs (a Batxillerat no se'n posen) ----------
\@ifundefined{estrelles}{}{%
  \renewcommand{\estrelles}[1]{{\color{curs}\ifcase#1\or\faStar\faStar[regular]\faStar[regular]\or\faStar\faStar\faStar[regular]\or\faStar\faStar\faStar\fi}}}

% ---------- capçalera i peu de les pàgines ----------
% primera pàgina: logos institucionals al peu
\fancypagestyle{primera}{\fancyhf{}\renewcommand{\headrulewidth}{0pt}
  \fancyfoot[L]{\raisebox{-0.32\height}{\includegraphics[height=0.78cm]{logo-gva.png}}\hspace{5mm}%
    \raisebox{-0.32\height}{\includegraphics[height=0.78cm]{logo-vermellar.png}}}
  \fancyfoot[R]{\scriptsize\color{muted}\href{https://creativecommons.org/licenses/by-nc-sa/4.0/deed.ca}{\ccbyncsa}\ \ CC BY-NC-SA 4.0 · fentmates.org\qquad\small\thepage}}
% la resta: el símbol a l'esquerra i el curs a la dreta, amb el títol al centre com abans
\fancyhead[L]{\raisebox{-0.3\height}{\includegraphics[height=0.62cm]{\fm@marca}}}
\fancyhead[R]{\small\bfseries\color{cursfosc}\fm@cursnom}
\renewcommand{\headrule}{\color{curs!45}\hrule height \headrulewidth}
\makeatother

\fmcurs{1batx-mi}

\begin{document}
\thispagestyle{primera}

\capcalera{Matemáticas I}{1\textsuperscript{o} Bach.}{Soluciones de la ficha · Unidad 5}{\raggedright Trigonometría}
\vspace{8pt}

\begin{tcolorbox}[colback=soft,colframe=soft,arc=2mm,left=3mm,right=3mm,top=1.5mm,bottom=1.5mm]
\small Intenta resolver cada ejercicio antes de mirar su solución. Los resultados se dan de manera exacta cuando es posible y, si no, con dos decimales; los ángulos, en grados, minutos y segundos. Puedes comprobar los triángulos, las razones y las ecuaciones con las herramientas de los apuntes interactivos de la unidad.
\end{tcolorbox}

% =====================================================
\seccio{1}{Medida de ángulos}

\exf{1} \begin{tasks}(4)
  \task $\res{\frac{5\pi}{4}}$
  \task $\res{\frac{5\pi}{12}}$
  \task $\res{\frac{8\pi}{9}}$
  \task $\res{-85^\circ56'37''}$
  \task $\res{114^\circ35'30''}$
  \task $\res{3\pi}$
  \task $\res{-\frac{3\pi}{2}}$
  \task $\res{\frac{211\pi}{270}}$$\approx2{,}4551$
  \task $\res{60^\circ}$
  \task $\res{17^\circ11'19''}$
  \task $\res{-216^\circ}$
  \task $\res{\frac{7\pi}{4}}$
\end{tasks}


% =====================================================
\seccio{2}{Trigonometría con la calculadora}

\exf{1} \begin{tasks}(3)
  \task $\sin(319^\circ12'52'')\approx\res{-0{,}6532}$
  \task $\cos(434^\circ26')\approx\res{0{,}2684}$
  \task $\text{tg}\,\frac{11\pi}{8}\approx\res{2{,}4142}$
  \task $\sec135^\circ\approx\res{-1{,}4142}$
  \task $\text{tg}\,7{,}03\approx\res{0{,}9257}$
  \task $\text{cosec}\,5{,}24\approx\res{-1{,}1574}$
  \task $\sin\frac{8\pi}{5}\approx\res{-0{,}9511}$
  \task $\text{cosec}\,(55'45'')\approx\res{61{,}6663}$
  \task $\cos\frac{\pi}{3}\approx\res{0{,}5000}$
  \task $\text{cosec}\,(200^\circ16')\approx\res{-2{,}8869}$
  \task $\sec\frac{5\pi}{4}\approx\res{-1{,}4142}$
  \task $\text{tg}\,(183^\circ13'53'')\approx\res{0{,}0565}$
  \task $\sin(333^\circ0'55'')\approx\res{-0{,}4538}$
  \task $\cos3{,}845\approx\res{-0{,}7626}$
  \task $\text{cotg}\,\frac{11\pi}{6}\approx\res{-1{,}7321}$
\end{tasks}

\ex{1} \begin{tasks}(2)
  \task $\cos x$: $\res{63^\circ0'21'',\ 296^\circ59'39''}$
  \task $\sin x$: $\res{68^\circ11'12'',\ 111^\circ48'48''}$
  \task $\text{tg}\,x$: $\res{151^\circ22'12'',\ 331^\circ22'12''}$
  \task $\cos x$: $\res{113^\circ33'11'',\ 246^\circ26'49''}$
  \task $\text{cotg}\,x$: $\res{24^\circ49'27'',\ 204^\circ49'27''}$
  \task $\cos x$: $\res{107^\circ0'49'',\ 252^\circ59'11''}$
  \task $\sin x$: $\res{43^\circ19'48'',\ 136^\circ40'12''}$
  \task $\sin x$: $\res{199^\circ19'2'',\ 340^\circ40'58''}$
\end{tasks}


% =====================================================
\seccio{3}{Razones trigonométricas de un ángulo agudo}

\exf{1} \begin{center}\renewcommand{\arraystretch}{1.5}\begin{tabular}{|l|*{5}{C{1.55cm}|}}\hline
\textbf{Grados} & $0^\circ$ & $30^\circ$ & $45^\circ$ & $60^\circ$ & $90^\circ$\\\hline
\textbf{Radianes} & $0$ & $\frac{\pi}{6}$ & $\frac{\pi}{4}$ & $\frac{\pi}{3}$ & $\frac{\pi}{2}$\\\hline
$\sin$ & $0$ & $\frac12$ & $\frac{\sqrt2}{2}$ & $\frac{\sqrt3}{2}$ & $1$\\\hline $\cos$ & $1$ & $\frac{\sqrt3}{2}$ & $\frac{\sqrt2}{2}$ & $\frac12$ & $0$\\\hline $\text{tg}$ & $0$ & $\frac{\sqrt3}{3}$ & $1$ & $\sqrt3$ & $\nexists$\\\hline $\text{cosec}$ & $\nexists$ & $2$ & $\sqrt2$ & $\frac{2\sqrt3}{3}$ & $1$\\\hline $\sec$ & $1$ & $\frac{2\sqrt3}{3}$ & $\sqrt2$ & $2$ & $\nexists$\\\hline $\text{cotg}$ & $\nexists$ & $\sqrt3$ & $1$ & $\frac{\sqrt3}{3}$ & $0$\\\hline
\end{tabular}\end{center}

\ex{1} \begin{tasks}(1)
  \task $\alpha$: $\sin\alpha=\res{\frac23}$, $\text{tg}\,\alpha=\res{\frac{2\sqrt5}{5}}$; $\sec\alpha=\frac{3\sqrt5}{5}$, $\text{cosec}\,\alpha=\frac32$, $\text{cotg}\,\alpha=\frac{\sqrt5}{2}$
  \task $\beta$: $\cos\beta=\res{\frac45}$, $\text{tg}\,\beta=\res{\frac34}$; $\sec\beta=\frac54$, $\text{cosec}\,\beta=\frac53$, $\text{cotg}\,\beta=\frac43$
  \task $\gamma$: $\cos\gamma=\res{\frac{\sqrt{10}}{10}}$, $\sin\gamma=\res{\frac{3\sqrt{10}}{10}}$; $\sec\gamma=\sqrt{10}$, $\text{cosec}\,\gamma=\frac{\sqrt{10}}{3}$, $\text{cotg}\,\gamma=\frac13$
  \task $\alpha$: $\sin\approx\res{0{,}5103}$, $\cos\approx\res{0{,}8600}$, $\text{tg}\approx\res{0{,}5934}$
  \task $\beta$: $\sin\approx\res{0{,}9429}$, $\cos\approx\res{0{,}3332}$, $\text{tg}\approx\res{2{,}8300}$
  \task $\gamma=45^\circ$: $\sin\gamma=\cos\gamma=\res{\frac{\sqrt2}{2}}$
\end{tasks}


% =====================================================
\seccio{4}{Razones trigonométricas de ángulos cualesquiera}

\exf{1} \begin{center}\renewcommand{\arraystretch}{1.45}\begin{tabular}{|C{1.4cm}|C{1.6cm}|C{1.6cm}|C{1.6cm}|C{1.6cm}|}\hline\textbf{Grados} & \textbf{Radianes} & $\sin$ & $\cos$ & $\text{tg}$\\\hline$120^\circ$ & $\frac{2\pi}{3}$ & $\frac{\sqrt3}{2}$ & $-\frac12$ & $-\sqrt3$\\\hline $135^\circ$ & $\frac{3\pi}{4}$ & $\frac{\sqrt2}{2}$ & $-\frac{\sqrt2}{2}$ & $-1$\\\hline $150^\circ$ & $\frac{5\pi}{6}$ & $\frac12$ & $-\frac{\sqrt3}{2}$ & $-\frac{\sqrt3}{3}$\\\hline $180^\circ$ & $\pi$ & $0$ & $-1$ & $0$\\\hline $210^\circ$ & $\frac{7\pi}{6}$ & $-\frac12$ & $-\frac{\sqrt3}{2}$ & $\frac{\sqrt3}{3}$\\\hline $225^\circ$ & $\frac{5\pi}{4}$ & $-\frac{\sqrt2}{2}$ & $-\frac{\sqrt2}{2}$ & $1$\\\hline $240^\circ$ & $\frac{4\pi}{3}$ & $-\frac{\sqrt3}{2}$ & $-\frac12$ & $\sqrt3$\\\hline $270^\circ$ & $\frac{3\pi}{2}$ & $-1$ & $0$ & $\nexists$\\\hline $300^\circ$ & $\frac{5\pi}{3}$ & $-\frac{\sqrt3}{2}$ & $\frac12$ & $-\sqrt3$\\\hline $315^\circ$ & $\frac{7\pi}{4}$ & $-\frac{\sqrt2}{2}$ & $\frac{\sqrt2}{2}$ & $-1$\\\hline $330^\circ$ & $\frac{11\pi}{6}$ & $-\frac12$ & $\frac{\sqrt3}{2}$ & $-\frac{\sqrt3}{3}$\\\hline $360^\circ$ & $2\pi$ & $0$ & $1$ & $0$\\\hline \end{tabular}\end{center}

\ex{1} \begin{tasks}(4)
  \task $\res{-2}$
  \task $\res{1}$
  \task $\res{-1}$
  \task $\res{1}$
\end{tasks}

\ex{2} \begin{center}\renewcommand{\arraystretch}{1.3}\begin{tabular}{|l|C{2.4cm}|C{2.4cm}|C{2.4cm}|}\hline\textbf{Cuadrante} & $\sin$ & $\cos$ & $\text{tg}$\\\hline Segundo & $0{,}6702$ & $\res{-0{,}7422}$ & $\res{-0{,}9030}$\\\hline Tercero & $\res{-0{,}8911}$ & $-0{,}4539$ & $\res{1{,}9631}$\\\hline Cuarto & $\res{-0{,}5979}$ & $\res{0{,}8016}$ & $-0{,}7459$\\\hline Tercero & $-0{,}7822$ & $\res{-0{,}6230}$ & $\res{1{,}2555}$\\\hline Segundo & $\res{0{,}8850}$ & $\res{-0{,}4657}$ & $-1{,}9004$\\\hline Cuarto & $\res{-0{,}7158}$ & $0{,}6983$ & $\res{-1{,}0251}$\\\hline \end{tabular}\end{center}

\ex{1} \begin{tasks}(1)
  \task $\cos\alpha=\frac14$, $\sin\alpha=\res{-\frac{\sqrt{15}}{4}}$ (negativo en el cuarto cuadrante), $\text{tg}\,\alpha=\res{-\sqrt{15}}$
  \task $\text{tg}\,\alpha=-2$, $\cos\alpha=-\frac{1}{\sqrt{1+4}}=\res{-\frac{\sqrt5}{5}}$, $\sin\alpha=\text{tg}\,\alpha\cdot\cos\alpha=\res{\frac{2\sqrt5}{5}}$
\end{tasks}

\ex{1} $\sin\alpha\approx\res{0{,}62}$, $\cos\alpha\approx\res{-0{,}78}$, $\text{tg}\,\alpha\approx\res{-0{,}79}$, $\text{cosec}\,\alpha\approx\res{1{,}61}$, $\sec\alpha\approx\res{-1{,}27}$, $\text{cotg}\,\alpha\approx\res{-1{,}27}$

\ex{1} $\sin\alpha\approx\res{-0{,}56}$, $\cos\alpha\approx\res{-0{,}83}$, $\text{tg}\,\alpha\approx\res{0{,}67}$, $\text{cosec}\,\alpha\approx\res{-1{,}79}$, $\sec\alpha\approx\res{-1{,}20}$, $\text{cotg}\,\alpha\approx\res{1{,}49}$

\ex{1} $\sin\alpha\approx\res{-0{,}68}$, $\cos\alpha\approx\res{0{,}74}$, $\text{tg}\,\alpha\approx\res{-0{,}92}$, $\text{cosec}\,\alpha\approx\res{-1{,}48}$, $\sec\alpha\approx\res{1{,}36}$, $\text{cotg}\,\alpha\approx\res{-1{,}09}$

\ex{1} \begin{tasks}(2)
  \task $\sin13290^\circ=\sin330^\circ=\res{-\frac12}$
  \task $\cos(-1680^\circ)=\cos120^\circ=\res{-\frac12}$
  \task $\text{tg}\,3825^\circ=\text{tg}\,225^\circ=\res{1}$
  \task $\cos4995^\circ=\cos315^\circ=\res{\frac{\sqrt2}{2}}$
  \task $\sin(-1710^\circ)=\sin90^\circ=\res{1}$
  \task $\text{tg}\,3630^\circ=\text{tg}\,30^\circ=\res{\frac{\sqrt3}{3}}$
  \task $\cos(-36000^\circ)=\cos0^\circ=\res{1}$
  \task $\sin(-330^\circ)=\sin30^\circ=\res{\frac12}$
\end{tasks}


% =====================================================
\seccio{5}{Relaciones entre las razones trigonométricas de algunos ángulos}

\exf{2} \begin{tasks}(1)
  \task $\cos\alpha=\res{-\frac23}$, $\sin\alpha=\res{\frac{\sqrt5}{3}}$, $\sin(180^\circ-\alpha)=\sin\alpha=\res{\frac{\sqrt5}{3}}$, $\cos(90^\circ+\alpha)=-\sin\alpha=\res{-\frac{\sqrt5}{3}}$, $\text{tg}\,(180^\circ+\alpha)=\text{tg}\,\alpha=\res{-\frac{\sqrt5}{2}}$
  \task $\alpha=180^\circ-\text{tg}^{-1}\frac{\sqrt5}{2}\approx\res{131^\circ48'37''}$
\end{tasks}

\ex{1} \begin{tasks}(3)
  \task $\res{-0{,}8}$
  \task $\res{0{,}6}$
  \task $\res{-0{,}75}$
  \task $\res{0{,}8}$
  \task $\res{-0{,}6}$
  \task $\res{\frac43}$
\end{tasks}

\ex{1} \begin{tasks}(2)
  \task $\cos48^\circ=\sin42^\circ=\res{0{,}67}$
  \task $\sin(-48^\circ)=-\cos42^\circ=\res{-0{,}74}$
  \task $\sin138^\circ=\sin42^\circ=\res{0{,}67}$
  \task $\text{tg}\,318^\circ=-\text{tg}\,42^\circ=\res{-0{,}9}$
  \task $\cos222^\circ=-\cos42^\circ=\res{-0{,}74}$
  \task $\text{tg}\,858^\circ=\text{tg}\,138^\circ=-\text{tg}\,42^\circ=\res{-0{,}9}$
\end{tasks}


% =====================================================
\seccio{6}{Resolución de triángulos rectángulos}

\exf{1} \begin{tasks}(1)
  \task $a=c\cos\hat B\approx\res{17{,}43\ \text{cm}}$, \ $b=c\sin\hat B\approx\res{26{,}84\ \text{cm}}$ \ ($\hat A=33^\circ$)
  \task $c=\sqrt{250^2+308^2}\approx\res{396{,}69\ \text{m}}$, \ $\text{tg}\,\hat A=\frac{250}{308}\Rightarrow\hat A\approx\res{39^\circ3'57''}$
  \task $b=\dfrac{a}{\text{tg}\,\hat A}\approx\res{56{,}01\ \text{cm}}$, \ $c=\dfrac{a}{\sin\hat A}\approx\res{66{,}05\ \text{cm}}$ \ ($\hat B=58^\circ$)
\end{tasks}

\ex{1} Altura: $12\sin50^\circ\approx\res{9{,}19\ \text{m}}$. Distancia: $12\cos50^\circ\approx\res{7{,}71\ \text{m}}$.

\ex{2} Con $A$ y $B$ a la misma altura, como en la figura, sumamos las proyecciones horizontales de los cuatro tramos: $4\cos60^\circ+1{,}7\cos30^\circ+3{,}5\cos20^\circ+4{,}2\cos50^\circ\approx2+1{,}47+3{,}29+2{,}7\approx\res{9{,}46\ \text{m}}$

\ex{2} Altura a los 6 km: $6\,\text{tg}\,15^\circ\approx1{,}608$ km. A los 11 km: $1{,}608+5\,\text{tg}\,25^\circ\approx\res{3{,}939\ \text{km}}>2{,}5$ km: \rest{pasa por encima de la montaña}.

\ex{2} La altura lo divide en dos triángulos rectángulos con un cateto de 5 m y un ángulo de $20^\circ$: lado igual $=\dfrac{5}{\sin20^\circ}\approx14{,}62$ m, altura $=\dfrac{5}{\text{tg}\,20^\circ}\approx13{,}74$ m. Perímetro $\approx\res{39{,}24\ \text{m}}$, área $\approx\res{68{,}69\ \text{m}^2}$.

\ex{1} $S=9\cdot16\cdot\sin65^\circ\approx\res{130{,}51\ \text{cm}^2}$


% =====================================================
\seccio{7}{Teorema del seno y teorema del coseno}

\exf{2} \begin{tasks}(1)
  \task Teorema del coseno: $\hat A\approx\res{44^\circ24'55''}$, $\hat B\approx\res{101^\circ32'13''}$, $\hat C\approx\res{34^\circ2'52''}$
  \task Teorema del coseno: $a\approx\res{5{,}77\ \text{cm}}$; $\hat B\approx\res{41^\circ4'43''}$, $\hat C\approx\res{99^\circ43'17''}$
  \task $\hat A=\res{74^\circ31'}$; teorema del seno: $b\approx\res{4{,}55\ \text{cm}}$, $c\approx\res{6{,}67\ \text{cm}}$
  \task Teorema del seno: $\sin\hat A\approx0{,}4774\Rightarrow\hat A\approx\res{28^\circ30'46''}$ (el obtuso no cabe); $\hat C\approx\res{114^\circ51'14''}$, $c\approx\res{15{,}21\ \text{cm}}$
  \task Teorema del coseno: $b\approx\res{11{,}27\ \text{cm}}$; $\hat A\approx\res{50^\circ58'46''}$, $\hat C\approx\res{25^\circ34'14''}$
\end{tasks}

\ex{1} \begin{tasks}(2)
  \task $12^2=144$ < $9^2+11^2=202$: \rest{acutángulo}
  \task \rest{no es un triángulo} ($23\geq8+14$)
  \task $26^2=676$ = $10^2+24^2=676$: \rest{rectángulo}
  \task $40^2=1600$ > $20^2+30^2=1300$: \rest{obtusángulo}
\end{tasks}


% =====================================================
\seccio{8}{Resolución de triángulos: problemas}

\exf{1} $8\sin38^\circ\approx\res{4{,}93\ \text{m}}$

\ex{1} $40\,\text{tg}\,26^\circ+1{,}75\approx19{,}51+1{,}75\approx\res{21{,}26\ \text{m}}$

\ex{2} Teorema del coseno: $\cos\theta=\dfrac{6^2+12^2-8^2}{2\cdot6\cdot12}=\dfrac{116}{144}\Rightarrow\theta\approx\res{36^\circ20'10''}$

\ex{2} $h=x\,\text{tg}\,40^\circ=(x+30)\,\text{tg}\,28^\circ\Rightarrow x=\dfrac{30\,\text{tg}\,28^\circ}{\text{tg}\,40^\circ-\text{tg}\,28^\circ}\approx\res{51{,}89\ \text{m}}$, \ $h\approx\res{43{,}54\ \text{m}}$

\ex{3} La visual al árbol cercano forma con la vertical un ángulo $\alpha$ con $\text{tg}\,\alpha=\frac{50}{25}=2$, es decir, $\alpha\approx63^\circ26'6''$. La del árbol lejano forma $\alpha+15^\circ$: $50+x=25\,\text{tg}\,(\alpha+15^\circ)\Rightarrow x\approx\res{72{,}17\ \text{m}}$

\ex{1} $d=\dfrac{40}{\sin15^\circ}\approx\res{154{,}55\ \text{m}}$

\ex{2} $\dfrac{h}{\text{tg}\,43^\circ}+\dfrac{h}{\text{tg}\,57^\circ50'}=15\Rightarrow h\approx\res{8{,}82\ \text{m}}$

\ex{2} \begin{tasks}(1)
  \task $\res{30^\circ,\ 60^\circ\ \text{y}\ 90^\circ}$
  \task Catetos $x$ y $x\sqrt3$: $\frac{x^2\sqrt3}{2}=28\Rightarrow x\approx\res{5{,}69}$, el otro cateto $\approx\res{9{,}85}$, hipotenusa $\approx\res{11{,}37}$ cm; perímetro $\approx\res{26{,}91\ \text{cm}}$
\end{tasks}

\ex{1} $\text{tg}\,\alpha=\frac{1{,}70}{0{,}50}=3{,}4\Rightarrow\alpha\approx\res{73^\circ36'38''}$

\ex{2} La anchura del tejado es $\dfrac{8}{\cos18^\circ}\approx8{,}41$ m: área $\approx12\cdot8{,}41\approx\res{100{,}94\ \text{m}^2}$

\ex{2} Después de 2 h han recorrido 10 y 8 km: $d=\sqrt{10^2+8^2-2\cdot10\cdot8\cos54^\circ}\approx\res{8{,}36\ \text{km}}$. Después de 6 h (el triple): $\approx\res{25{,}09\ \text{km}}$.

\ex{3} \begin{tasks}(1)
  \task El tercer ángulo es $93^\circ$. Teorema del seno: tramos de $\dfrac{8\sin37^\circ}{\sin93^\circ}\approx4{,}82$ m y $\dfrac{8\sin50^\circ}{\sin93^\circ}\approx6{,}14$ m: cuerda $\approx\res{10{,}96\ \text{m}}$
  \task $d=4{,}82\cdot\sin50^\circ\approx\res{3{,}69\ \text{m}}$
\end{tasks}

\ex{3} Altura $h=8\sin65^\circ\approx7{,}25$; proyección del lado $8\cos65^\circ\approx3{,}38$. Por Pitágoras con la diagonal, la base mayor es $B\approx13{,}13+3{,}38\approx16{,}51$ y la menor $b\approx9{,}75$ cm. Área $=\frac{B+b}{2}\,h\approx\res{95{,}21\ \text{cm}^2}$

\ex{2} Proyección del lado: $\frac{14-8}{2}=3$ cm: $\cos\alpha=\frac35\Rightarrow\alpha\approx\res{53^\circ7'48''}$ (dos ángulos) y $180^\circ-\alpha\approx\res{126^\circ52'12''}$ (los otros dos)

\ex{1} $2\,\text{tg}^{-1}\frac{6}{10}\approx\res{61^\circ55'39''}$ \ (y el suplementario, $\res{118^\circ4'21''}$)

\ex{2} $\cos\theta=\dfrac{10^2+10^2-15^2}{2\cdot10\cdot10}=-0{,}125\Rightarrow\theta\approx\res{97^\circ10'51''}$; los otros dos: $\res{41^\circ24'35''}$

\ex{2} $\cos\alpha=\dfrac{5^2+8^2-10^2}{2\cdot5\cdot8}=-\frac{11}{80}\Rightarrow\alpha\approx\res{97^\circ54'12''}$ (dos ángulos) y $\res{82^\circ5'48''}$ (los otros dos)


% =====================================================
\seccio{9}{Fórmulas trigonométricas}

\exf{1} $15^\circ=45^\circ-30^\circ$: $\sin=\res{\frac{\sqrt6-\sqrt2}{4}}$, $\cos=\res{\frac{\sqrt6+\sqrt2}{4}}$, $\text{tg}=\res{2-\sqrt3}$. \quad $\frac{5\pi}{12}=75^\circ=45^\circ+30^\circ$: $\sin=\res{\frac{\sqrt6+\sqrt2}{4}}$, $\cos=\res{\frac{\sqrt6-\sqrt2}{4}}$, $\text{tg}=\res{2+\sqrt3}$.

\ex{1} $\cos\alpha=0{,}8$, $\cos\beta=0{,}6$.\begin{tasks}(3)
  \task $\res{1}$
  \task $\res{-0{,}28}$
  \task $\res{0}$
  \task $\res{0{,}96}$
  \task $\res{\nexists}$ ($\cos=0$)
  \task $\res{-\frac{7}{24}}$
\end{tasks}

\ex{1} $\sin\alpha\cos\frac{3\pi}{2}+\cos\alpha\sin\frac{3\pi}{2}=\sin\alpha\cdot0+\cos\alpha\cdot(-1)=\res{-\cos\alpha}$

\ex{2} $\sin10^\circ=\sqrt{\frac{1-0{,}94}{2}}\approx\res{0{,}1732}$, $\cos10^\circ=\sqrt{\frac{1+0{,}94}{2}}\approx\res{0{,}9849}$; \ $\sin40^\circ=2\cdot0{,}34\cdot0{,}94\approx\res{0{,}6392}$, $\cos40^\circ=0{,}94^2-0{,}34^2=\res{0{,}768}$

\ex{2} $\cos\alpha=-\frac45$; $45^\circ<\frac\alpha2<90^\circ$: $\sin\frac\alpha2=\res{\frac{3\sqrt{10}}{10}}$, $\cos\frac\alpha2=\res{\frac{\sqrt{10}}{10}}$, $\text{tg}\,\frac\alpha2=\res{3}$

\ex{3} $\cos\alpha=-\frac45$, $\text{tg}\,\alpha=-\frac34$, $\sin\beta=-\frac{\sqrt{15}}{4}$. \ $\cos(\alpha+\beta)=\res{\frac{4+3\sqrt{15}}{20}}$, \ $\sin(\alpha-\beta)=\res{-\frac{3+4\sqrt{15}}{20}}$, \ $\text{tg}\,2\alpha=\res{-\frac{24}{7}}$, \ $\text{tg}\,\frac\beta2=\res{-\frac{\sqrt{15}}{3}}$ ($\frac\beta2$ es del segundo cuadrante)

\ex{2} En los dos casos: $\sin2\alpha=\res{\frac35}$, $\cos2\alpha=\res{-\frac45}$, $\text{tg}\,2\alpha=\res{-\frac34}$ (los dos valores de $2\alpha$ difieren en $360^\circ$).

\ex{2} $45^\circ<\frac\alpha2<90^\circ$: $\sin\frac\alpha2=\sqrt{\frac56}=\res{\frac{\sqrt{30}}{6}}$, $\cos\frac\alpha2=\sqrt{\frac16}=\res{\frac{\sqrt6}{6}}$, $\text{tg}\,\frac\alpha2=\res{\sqrt5}$

\ex{3} \begin{tasks}(1)
  \task $\dfrac{\sin\alpha-\cos\alpha}{\frac{\sin\alpha-\cos\alpha}{\cos\alpha}}=\cos\alpha$
  \task $\dfrac{\sin^2\alpha-\sin^2\alpha\cos^2\alpha}{\cos^2\alpha}=\dfrac{\sin^2\alpha(1-\cos^2\alpha)}{\cos^2\alpha}=\text{tg}^2\alpha\sin^2\alpha$
  \task $\dfrac{\frac{\sin\alpha+\cos\alpha}{\sin\alpha}}{\sin\alpha+\cos\alpha}=\dfrac{1}{\sin\alpha}$
  \task $\dfrac{1-\cos^2\alpha}{\cos^2\alpha}=\dfrac{\sin^2\alpha}{\cos^2\alpha}$
  \task $\dfrac{\sin^2\alpha+\cos^2\alpha+2\sin\alpha\cos\alpha}{2\cos^2\alpha}=\dfrac{(\sin\alpha+\cos\alpha)^2}{2\cos^2\alpha}=\frac12(\text{tg}\,\alpha+1)^2$
  \task $2\,\text{tg}\,\alpha\cdot\dfrac{1-\cos\alpha}{2}+\sin\alpha=\text{tg}\,\alpha-\sin\alpha+\sin\alpha$
  \task $\cos2\alpha\cos\alpha-\sin2\alpha\sin\alpha=(2\cos^2\alpha-1)\cos\alpha-2(1-\cos^2\alpha)\cos\alpha$
  \task $\dfrac{\sin2\alpha}{\cos2\alpha}-\dfrac{\sin\alpha}{\cos\alpha}=\dfrac{\sin(2\alpha-\alpha)}{\cos2\alpha\cos\alpha}=\dfrac{\text{tg}\,\alpha}{\cos2\alpha}$
  \task $\dfrac{\sin^2\alpha+\cos^2\alpha}{\sin\alpha\cos\alpha}=\dfrac{1}{\cos\alpha}\cdot\dfrac{1}{\sin\alpha}$
  \task $\dfrac{2\sin\alpha\cos\alpha}{2\cos^2\alpha}=\dfrac{\sin\alpha}{\cos\alpha}$
  \task $\cos^2\alpha+\dfrac{\cos^4\alpha}{\sin^2\alpha}=\dfrac{\cos^2\alpha(\sin^2\alpha+\cos^2\alpha)}{\sin^2\alpha}=\dfrac{1}{\text{tg}^2\alpha}$
\end{tasks}

\ex{3} \begin{tasks}(1)
  \task Es $\cos\big(\alpha-(\alpha-\beta)\big)=\cos\beta$.
  \task $\sin^2\alpha\cos^2\beta-\cos^2\alpha\sin^2\beta=\sin^2\alpha(1-\sin^2\beta)-(1-\sin^2\alpha)\sin^2\beta$
  \task $\dfrac{2\cos\alpha\cos\beta}{2\sin\alpha\cos\beta}$
  \task $2-2(\cos\alpha\cos\beta-\sin\alpha\sin\beta)=2-2\cos(\alpha+\beta)=4\sin^2\frac{\alpha+\beta}{2}$
\end{tasks}


% =====================================================
\seccio{10}{Ecuaciones trigonométricas}

\exf{1} \begin{tasks}(1)
  \task $\sin x=\frac12$: $\res{30^\circ}$, $\res{150^\circ}$\ \ \small($x=30^\circ+360^\circ k,\ 150^\circ+360^\circ k$)
  \task $\cos x=\frac{\sqrt3}{2}$: $\res{30^\circ}$, $\res{330^\circ}$\ \ \small($x=\pm30^\circ+360^\circ k$)
  \task $\text{tg}\,x=1$: $\res{45^\circ}$, $\res{225^\circ}$\ \ \small($x=45^\circ+180^\circ k$)
  \task $\sin x=-\frac{\sqrt2}{2}$: $\res{225^\circ}$, $\res{315^\circ}$\ \ \small($x=225^\circ+360^\circ k,\ 315^\circ+360^\circ k$)
  \task $\cos x=-\frac12$: $\res{120^\circ}$, $\res{240^\circ}$\ \ \small($x=\pm120^\circ+360^\circ k$)
  \task $1+\cos x=0$: $\res{180^\circ}$\ \ \small($x=180^\circ+360^\circ k$)
  \task $\sin4x=-\frac{\sqrt3}{2}$: $\res{60^\circ}$, $\res{75^\circ}$, $\res{150^\circ}$, $\res{165^\circ}$, $\res{240^\circ}$, $\res{255^\circ}$, $\res{330^\circ}$, $\res{345^\circ}$\ \ \small($x=60^\circ+90^\circ k,\ 75^\circ+90^\circ k$)
  \task $\cos2x=\frac{\sqrt2}{2}$: $\res{22^\circ30'}$, $\res{157^\circ30'}$, $\res{202^\circ30'}$, $\res{337^\circ30'}$\ \ \small($x=\pm22^\circ30^{\prime}+180^\circ k$)
  \task $\text{tg}\,3x=-1$: $\res{45^\circ}$, $\res{105^\circ}$, $\res{165^\circ}$, $\res{225^\circ}$, $\res{285^\circ}$, $\res{345^\circ}$\ \ \small($x=45^\circ+60^\circ k$)
  \task $\sin\frac x2=0$: $\res{0^\circ}$\ \ \small($x=360^\circ k$)
  \task $\cos\frac x3=-\frac12$: \rest{ninguna} en $[0^\circ,360^\circ)$\ \ \small($x=360^\circ+1080^\circ k,\ 720^\circ+1080^\circ k$)
  \task $\text{tg}\,\frac{3x}{4}=-\frac{\sqrt3}{3}$: $\res{200^\circ}$\ \ \small($x=200^\circ+240^\circ k$)
\end{tasks}

\ex{2} \begin{tasks}(1)
  \task $\cos(2x+30^\circ)=\pm\frac{\sqrt2}{2}\Rightarrow2x+30^\circ=45^\circ+90^\circ k\Rightarrow x=7^\circ30^{\prime}+45^\circ k$: $\res{7^\circ30'}$, $\res{52^\circ30'}$, $\res{97^\circ30'}$, $\res{142^\circ30'}$, $\res{187^\circ30'}$, $\res{232^\circ30'}$, $\res{277^\circ30'}$, $\res{322^\circ30'}$
  \task $\cos x(\cos x-\sqrt3\sin x)=0$: $\cos x=0$ o $\text{tg}\,x=\frac{\sqrt3}{3}$: $\res{30^\circ}$, $\res{90^\circ}$, $\res{210^\circ}$, $\res{270^\circ}$
  \task $2\cos(x+30^\circ)=1\Rightarrow x+30^\circ=60^\circ$ o $300^\circ$: $\res{30^\circ}$, $\res{270^\circ}$
  \task $\sin^2x=-2\cos x\Rightarrow\cos^2x-2\cos x-1=0\Rightarrow\cos x=1-\sqrt2$: $\res{114^\circ28'11''}$, $\res{245^\circ31'49''}$
  \task $\dfrac{2\sin2x\cos x}{2\cos2x\cos x}=\text{tg}\,2x=\sqrt3$: $\res{30^\circ}$, $\res{120^\circ}$, $\res{210^\circ}$, $\res{300^\circ}$
  \task $\sin x(\sin^2x-\cos^2x)=-\sin x\cos2x=0$: $\res{0^\circ}$, $\res{45^\circ}$, $\res{135^\circ}$, $\res{180^\circ}$, $\res{225^\circ}$, $\res{315^\circ}$
  \task $2\sin(x+30^\circ)=2\Rightarrow x+30^\circ=90^\circ$: $\res{60^\circ}$
  \task $\dfrac{1-\cos x}{1+\cos x}=1-\cos x\Rightarrow\cos x=1$ o $\cos x=0$: $\res{0^\circ}$, $\res{90^\circ}$, $\res{270^\circ}$
  \task $\dfrac{2\cos3x\cos x}{2\cos3x\sin x}=\text{cotg}\,x=1$: $\res{45^\circ}$, $\res{225^\circ}$
  \task Dividiendo entre $\sin x\neq0$: $4\cos^2x+4\cos x+1=0\Rightarrow\cos x=-\frac12$: $\res{120^\circ}$, $\res{240^\circ}$
\end{tasks}

\ex{2} \begin{tasks}(1)
  \task $\text{tg}^2x-4\,\text{tg}\,x+3=0\Rightarrow\text{tg}\,x=1$ o $3$: $\res{45^\circ}$, $\res{71^\circ33'54''}$, $\res{225^\circ}$, $\res{251^\circ33'54''}$
  \task $2\sin x(\cos^2x-3\sin^2x)=0$: $\sin x=0$ o $\text{tg}^2x=\frac13$: $\res{0^\circ}$, $\res{30^\circ}$, $\res{150^\circ}$, $\res{180^\circ}$, $\res{210^\circ}$, $\res{330^\circ}$
  \task $8\sin^2x+2\sin x-3=0\Rightarrow\sin x=\frac12$ o $-\frac34$: $\res{30^\circ}$, $\res{150^\circ}$, $\res{228^\circ35'25''}$, $\res{311^\circ24'35''}$
  \task $\text{tg}\,x=\frac{\sqrt3}{3}$: $\res{30^\circ}$, $\res{210^\circ}$
  \task $\sin x(2\cos^2x-1)=0$: $\res{0^\circ}$, $\res{45^\circ}$, $\res{135^\circ}$, $\res{180^\circ}$, $\res{225^\circ}$, $\res{315^\circ}$
  \task $2\cos2x\cos x=0$: $\res{45^\circ}$, $\res{90^\circ}$, $\res{135^\circ}$, $\res{225^\circ}$, $\res{270^\circ}$, $\res{315^\circ}$
  \task $2\cos^2x-\cos x-3=0\Rightarrow\cos x=-1$ ($\frac32$ no vale): $\res{180^\circ}$
  \task $\sqrt2\sin^2x+\sin x-\sqrt2=0\Rightarrow\sin x=\frac{\sqrt2}{2}$: $\res{45^\circ}$, $\res{135^\circ}$
  \task $2\cos2x=7\sin x\Rightarrow4\sin^2x+7\sin x-2=0\Rightarrow\sin x=\frac14$ ($-2$ no vale): $\res{14^\circ28'39''}$, $\res{165^\circ31'21''}$
  \task $4\cos^2x+8\cos x-5=0\Rightarrow\cos x=\frac12$ ($-\frac52$ no vale): $\res{60^\circ}$, $\res{300^\circ}$
  \task $-2\sin^2x-4\sin x=0\Rightarrow\sin x=0$ ($-2$ no vale): $\res{0^\circ}$, $\res{180^\circ}$
\end{tasks}


% =====================================================
\seccio{11}{Sistemas de ecuaciones trigonométricas}

\exf{2} \begin{tasks}(1)
  \task $\cos y=1+\sin x\Rightarrow(2\sin x+1)^2=0$: $\sin x=-\frac12$, $\cos y=\frac12$: $\res{(210^\circ,60^\circ)}$, $\res{(210^\circ,300^\circ)}$, $\res{(330^\circ,60^\circ)}$, $\res{(330^\circ,300^\circ)}$
  \task $\cos x$ y $\cos y$ son las raíces de $t^2+\frac12t-\frac12=0$: $\frac12$ y $-1$: $\res{(60^\circ,180^\circ)}$, $\res{(180^\circ,60^\circ)}$, $\res{(180^\circ,300^\circ)}$, $\res{(300^\circ,180^\circ)}$
  \task $\sin x\sin y=\frac{(\sin x+\sin y)^2-(\sin^2x+\sin^2y)}{2}=0$: un seno vale 0 y el otro $\frac{\sqrt3}{2}$: $\res{(0^\circ,60^\circ)}$, $\res{(0^\circ,120^\circ)}$, $\res{(60^\circ,0^\circ)}$, $\res{(60^\circ,180^\circ)}$, $\res{(120^\circ,0^\circ)}$, $\res{(120^\circ,180^\circ)}$, $\res{(180^\circ,60^\circ)}$, $\res{(180^\circ,120^\circ)}$
  \task $\cos y=\sin x\Rightarrow\sqrt3\cos x-\sin x=1\Rightarrow\cos(x+30^\circ)=\frac12$: $\res{x=\frac{\pi}{6},\ y=\frac{\pi}{3}}$ (con $x=\frac{3\pi}{2}$, $y=-\pi$ quedaría fuera del intervalo)
  \task Sumando y restando: $\sin(x+y)=\frac12$, $\sin(x-y)=0$: $\res{(15^\circ,15^\circ)}$, $\res{(75^\circ,75^\circ)}$, $\res{(105^\circ,285^\circ)}$, $\res{(165^\circ,345^\circ)}$, $\res{(195^\circ,195^\circ)}$, $\res{(255^\circ,255^\circ)}$, $\res{(285^\circ,105^\circ)}$, $\res{(345^\circ,165^\circ)}$
\end{tasks}

\ex{3} \begin{tasks}(1)
  \task $y=\frac\pi2-x\Rightarrow\sin x+\cos x=1$: $\res{(0,\tfrac\pi2)}$ y $\res{(\tfrac\pi2,0)}$
  \task $\sin\left(y+\frac\pi3\right)=2\sin y\Rightarrow\text{tg}\,y=\frac{\sqrt3}{3}$: $\res{y=\frac\pi6,\ x=\frac\pi2}$ y $\res{y=\frac{7\pi}{6},\ x=\frac{3\pi}{2}}$
  \task $2y\sin2x=3$ y $2y\cos2x=\sqrt3$: dividiendo, $\text{tg}\,2x=\sqrt3$; sumando los cuadrados, $4y^2=12$. En $[0,2\pi)$: $\res{x=\frac\pi6,\frac{7\pi}{6}\ \text{con}\ y=\sqrt3}$ y $\res{x=\frac{2\pi}{3},\frac{5\pi}{3}\ \text{con}\ y=-\sqrt3}$
\end{tasks}


% =====================================================
\seccio{12}{Funciones trigonométricas}

\exf{1} \begin{tasks}(4)
  \task \rest{V}
  \task \rest{V}
  \task \rest{V}
  \task \rest{V} ($\cos x=\sin(x+\frac\pi2)$)
\end{tasks}

\ex{1} \begin{tasks}(1)
  \task $\frac x2=\frac\pi2+k\pi$: $\res{x=\pi,\ 3\pi}$
  \task $x-\pi=k\pi$: $\res{x=0,\ \pi,\ 2\pi,\ 3\pi,\ 4\pi}$
  \task $x+\pi=\frac\pi2+k\pi$: $\res{x=\frac\pi2,\ \frac{3\pi}2,\ \frac{5\pi}2,\ \frac{7\pi}2}$
\end{tasks}

\ex{2} \begin{tasks}(2)
  \task $=\cos x$: \rest{la gráfica del coseno}
  \task $=-\sin x$: \rest{la del seno reflejada en el eje $X$}
  \task $=\cos x$: \rest{la gráfica del coseno}
  \task $=\sin x$: \rest{la gráfica del seno}
\end{tasks}

\ex{2} \begin{tasks}(4)
  \task a) $y=2\cos x$ $\to$ \res{(3)}
  \task b) $y=\cos2x$ $\to$ \res{(7)}
  \task c) $y=\frac12\sin x$ $\to$ \res{(8)}
  \task d) $y=\cos\frac x2$ $\to$ \res{(2)}
  \task e) $y=\sin2x$ $\to$ \res{(6)}
  \task f) $y=2\sin x$ $\to$ \res{(1)}
  \task g) $y=\frac12\cos x$ $\to$ \res{(5)}
  \task h) $y=\sin\frac x2$ $\to$ \res{(4)}
\end{tasks}Fíjate en la amplitud (el valor máximo) y en el período: $2\pi$, $\pi$ (si hay $2x$) o $4\pi$ (si hay $\frac x2$).


% =====================================================
\seccio{13}{Ampliación y problemas (ejercicios nuevos)}

\exf{1} 25 min $=150^\circ=\frac{5\pi}{6}$ rad: $l=12\cdot\frac{5\pi}{6}=\res{10\pi\approx31{,}42\ \text{cm}}$; \ $S=\frac12\cdot144\cdot\frac{5\pi}{6}=\res{60\pi\approx188{,}5\ \text{cm}^2}$

\ex{1} $\frac{\sqrt3}{2}\cdot\frac{\sqrt2}{2}+\frac{\sqrt3}{3}=\res{\frac{\sqrt6}{4}+\frac{\sqrt3}{3}}$; \quad $\dfrac{\frac12+\frac12}{-1}=\res{-1}$

\ex{2} $\sin x-\sin x-\sin x=\res{-\sin x}$; \quad $\dfrac{(-\cos x)(-\text{tg}\,x)}{-\sin x}=\dfrac{\sin x}{-\sin x}=\res{-1}$

\ex{2} $\sin\hat B=\frac{8\sin40^\circ}{6}\approx0{,}8571$: \rest{dos soluciones}. \ $\hat B_1\approx\res{58^\circ59'13''}$, $\hat C_1\approx81^\circ0'47''$, $c_1\approx\res{9{,}22}$ cm; \ $\hat B_2\approx\res{121^\circ0'47''}$, $\hat C_2\approx18^\circ59'13''$, $c_2\approx\res{3{,}04}$ cm.

\ex{2} $\hat A\approx\res{53^\circ7'48''}$, $\hat B\approx\res{59^\circ29'23''}$, $\hat C\approx\res{67^\circ22'48''}$. \ $p=21$, $S=\sqrt{21\cdot8\cdot7\cdot6}=\res{84\ \text{cm}^2}$

\ex{2} Han recorrido 27 y 36 km: $d=\sqrt{27^2+36^2-2\cdot27\cdot36\cos70^\circ}\approx\res{36{,}88\ \text{km}}$

\ex{2} $h=x\,\text{tg}\,35^\circ=(x+20)\,\text{tg}\,22^\circ\Rightarrow$ anchura $x\approx\res{27{,}28\ \text{m}}$, altura $h\approx\res{19{,}1\ \text{m}}$

\ex{2} Cada triángulo central tiene un ángulo de $72^\circ$: lado $=2\cdot10\sin36^\circ\approx\res{11{,}76\ \text{cm}}$; área $=5\cdot\frac12\cdot10^2\sin72^\circ\approx\res{237{,}76\ \text{cm}^2}$

\ex{3} $\sin\alpha=-\frac{2\sqrt5}{5}$, $\cos\alpha=-\frac{\sqrt5}{5}$, $\sin\beta=\frac45$, $\text{tg}\,\beta=-\frac43$. \ $\sin(\alpha+\beta)=\res{\frac{2\sqrt5}{25}}$, $\cos(\alpha-\beta)=\res{-\frac{\sqrt5}{5}}$, $\text{tg}\,(\alpha+\beta)=\res{\frac{2}{11}}$, $\sin2\alpha=\res{\frac45}$

\ex{2} $(\cos^2\alpha+\sin^2\alpha)(\cos^2\alpha-\sin^2\alpha)=1\cdot\cos2\alpha$; \quad $\dfrac{2\sin^2\alpha}{2\sin\alpha\cos\alpha}=\text{tg}\,\alpha$

\ex{2} $2\cos^2x-3\cos x+1=0\Rightarrow\cos x=1$ o $\frac12$: $\res{0^\circ,\ 60^\circ,\ 300^\circ}$

\ex{2} $\cos x(2\sin x-1)=0$: $\res{30^\circ,\ 90^\circ,\ 150^\circ,\ 270^\circ}$

\ex{2} Amplitud $\res{2}$, período $\res{\pi}$, recorrido $\res{[-1,3]}$, desplazada $\res{\frac\pi4}$ hacia la derecha. Máximo ($y=3$) cuando $2x-\frac\pi2=\frac\pi2$: $\res{x=\frac\pi2}$


\end{document}
